% GENERATED FILE. Edit canonical lessons, then run npm run generate:books.
% canonical-source-hash: fnv1a64:b766f74b12b9429a
% canonical-lessons: KA-C260-samanu, KA-C260-kili, KA-C260-patrike, KA-C260-vahana, KA-C260-ato

\chapter{Luggage, Key, Newspaper, Vehicle, Autorickshaw}
\label{ch:ka-luggage-key-newspaper-vehicle-autorickshaw}
\hlchaptermodality{260}

\begin{chapteropening}
\textbf{By the end of this chapter:} \emph{I can say luggage, a key, a newspaper, a vehicle and an autorickshaw.}

Complete the last lesson of chapter 260: I can say luggage, a key, a newspaper, a vehicle and an autorickshaw.
\end{chapteropening}

\section[\texorpdfstring{\kn{ಸಾಮಾನು} (sāmānu)}{ಸಾಮಾನು (sāmānu)}]{\kn{ಸಾಮಾನು} (sāmānu) — luggage, belongings}

\label{lesson:KA-C260-samanu}



\begin{quote}
Before the new one: say the Kannada for a handkerchief, then the Kannada for a necklace.
\end{quote}

\subsection*{You'll want to know: \kn{ಸಾಮಾನು}}
\textbf{\kn{ಸಾಮಾನು}} — \emph{sāmānu} — \textquotedblleft{}luggage, belongings\textquotedblright{}.

\begin{cousinweb}
\textbf{\kn{ಸಾಮಾನು}} (\emph{sāmānu}) is Persian \emph{sāmān}, \textquotedblleft{}goods, equipment\textquotedblright{}, which came into Kannada through Hindustani, with a final \emph{-u} added.
\end{cousinweb}

\begin{grammarlens}[title={What you've built}]
One more word for this chapter.
\end{grammarlens}

\subsection*{Your turn}
\begin{itemize}
\raggedright
  \item \emph{Say it:} \emph{sāmānu}
  \item \emph{Say it:} \emph{sāmānu}, once more
  \item \emph{Say it:} say \emph{sāmānu}
  \item \emph{Recall:} say the Kannada for a handkerchief, then the Kannada for a necklace, then say \emph{sāmānu} again
\end{itemize}

\begin{tcolorbox}[breakable,colback=teal!4,colframe=teal!35!black,title={Before you move on}]
What does \kn{ಸಾಮಾನು} mean? (\textquotedblleft{}Luggage, belongings\textquotedblright{}.) Say it once more.
\end{tcolorbox}
\section[\texorpdfstring{\kn{ಕೀಲಿ} (kīli)}{ಕೀಲಿ (kīli)}]{\kn{ಕೀಲಿ} (kīli) — a key}

\label{lesson:KA-C260-kili}



\begin{quote}
Before the new one: say the Kannada for a necklace, then the Kannada for luggage.
\end{quote}

\subsection*{You'll want to know: \kn{ಕೀಲಿ}}
\textbf{\kn{ಕೀಲಿ}} — \emph{kīli} — \textquotedblleft{}a key\textquotedblright{}.

\begin{cousinweb}
\textbf{\kn{ಕೀಲಿ}} (\emph{kīli}) is probably the same word as \emph{kīlu}, \textquotedblleft{}a pin, a bolt\textquotedblright{}, which DEDR lists among Indo-Aryan loans (DEDR App. 26): it goes back to Sanskrit \emph{kīla}, \textquotedblleft{}a stake, a bolt\textquotedblright{} (CDIAL 3202). CDIAL 3202 adds that the Sanskrit word itself may be of non-Aryan origin.
\end{cousinweb}

\begin{grammarlens}[title={What you've built}]
One more word for this chapter.
\end{grammarlens}

\subsection*{Your turn}
\begin{itemize}
\raggedright
  \item \emph{Say it:} \emph{kīli}
  \item \emph{Say it:} \emph{kīli}, once more
  \item \emph{Say it:} say \emph{kīli}
  \item \emph{Recall:} say the Kannada for a necklace, then the Kannada for luggage, then say \emph{kīli} again
\end{itemize}

\begin{tcolorbox}[breakable,colback=teal!4,colframe=teal!35!black,title={Before you move on}]
What does \kn{ಕೀಲಿ} mean? (\textquotedblleft{}A key\textquotedblright{}.) Say it once more.
\end{tcolorbox}
\section[\texorpdfstring{\kn{ಪತ್ರಿಕೆ} (patrike)}{ಪತ್ರಿಕೆ (patrike)}]{\kn{ಪತ್ರಿಕೆ} (patrike) — a newspaper}

\label{lesson:KA-C260-patrike}



\begin{quote}
Before the new one: say the Kannada for luggage, then the Kannada for a key.
\end{quote}

\subsection*{You'll want to know: \kn{ಪತ್ರಿಕೆ}}
\textbf{\kn{ಪತ್ರಿಕೆ}} — \emph{patrike} — \textquotedblleft{}a newspaper\textquotedblright{}.

\begin{cousinweb}
\textbf{\kn{ಪತ್ರಿಕೆ}} (\emph{patrike}) is Sanskrit \emph{patrikā}, \textquotedblleft{}a leaf, a letter\textquotedblright{}, which CDIAL 7733 spells \emph{pattrikā}, with the Kannada \emph{-e} for final \emph{-ā}: a newspaper is \textquotedblleft{}a sheet\textquotedblright{}.
\end{cousinweb}

\begin{grammarlens}[title={What you've built}]
One more word for this chapter.
\end{grammarlens}

\subsection*{Your turn}
\begin{itemize}
\raggedright
  \item \emph{Say it:} \emph{patrike}
  \item \emph{Say it:} \emph{patrike}, once more
  \item \emph{Say it:} say \emph{patrike}
  \item \emph{Recall:} say the Kannada for luggage, then the Kannada for a key, then say \emph{patrike} again
\end{itemize}

\begin{tcolorbox}[breakable,colback=teal!4,colframe=teal!35!black,title={Before you move on}]
What does \kn{ಪತ್ರಿಕೆ} mean? (\textquotedblleft{}A newspaper\textquotedblright{}.) Say it once more.
\end{tcolorbox}
\section[\texorpdfstring{\kn{ವಾಹನ} (vāhana)}{ವಾಹನ (vāhana)}]{\kn{ವಾಹನ} (vāhana) — a vehicle}

\label{lesson:KA-C260-vahana}



\begin{quote}
Before the new one: say the Kannada for a key, then the Kannada for a newspaper.
\end{quote}

\subsection*{You'll want to know: \kn{ವಾಹನ}}
\textbf{\kn{ವಾಹನ}} — \emph{vāhana} — \textquotedblleft{}a vehicle\textquotedblright{}.

\begin{cousinweb}
\textbf{\kn{ವಾಹನ}} (\emph{vāhana}) is Sanskrit \emph{vāhana}, \textquotedblleft{}any vehicle or draught animal\textquotedblright{}, taken over whole as a learned loan, a tatsama (CDIAL 11610).
\end{cousinweb}

\begin{grammarlens}[title={What you've built}]
One more word for this chapter.
\end{grammarlens}

\subsection*{Your turn}
\begin{itemize}
\raggedright
  \item \emph{Say it:} \emph{vāhana}
  \item \emph{Say it:} \emph{vāhana}, once more
  \item \emph{Say it:} say \emph{vāhana}
  \item \emph{Recall:} say the Kannada for a key, then the Kannada for a newspaper, then say \emph{vāhana} again
\end{itemize}

\begin{tcolorbox}[breakable,colback=teal!4,colframe=teal!35!black,title={Before you move on}]
What does \kn{ವಾಹನ} mean? (\textquotedblleft{}A vehicle\textquotedblright{}.) Say it once more.
\end{tcolorbox}
\section[\texorpdfstring{\kn{ಆಟೋ} (āṭō)}{ಆಟೋ (āṭō)}]{\kn{ಆಟೋ} (āṭō) — an autorickshaw}

\label{lesson:KA-C260-ato}



\begin{quote}
Before the new one: say the Kannada for a newspaper, then the Kannada for a vehicle.
\end{quote}

\subsection*{You'll want to know: \kn{ಆಟೋ}}
\textbf{\kn{ಆಟೋ}} — \emph{āṭō} — \textquotedblleft{}an autorickshaw\textquotedblright{}.

\begin{cousinweb}
\textbf{\kn{ಆಟೋ}} (\emph{āṭō}) is a loan from English \emph{auto}. It is short for \textquotedblleft{}autorickshaw\textquotedblright{}.
\end{cousinweb}

\begin{grammarlens}[title={What you've built}]
One more word for this chapter.
\end{grammarlens}

\subsection*{Your turn}
\begin{itemize}
\raggedright
  \item \emph{Say it:} \emph{āṭō}
  \item \emph{Say it:} \emph{āṭō}, once more
  \item \emph{Say it:} say \emph{āṭō}
  \item \emph{Recall:} say the Kannada for a newspaper, then the Kannada for a vehicle, then say \emph{āṭō} again
\end{itemize}

\begin{tcolorbox}[breakable,colback=teal!4,colframe=teal!35!black,title={Before you move on}]
What does \kn{ಆಟೋ} mean? (\textquotedblleft{}An autorickshaw\textquotedblright{}.) Say it once more.
\end{tcolorbox}
